%-------------------------------------------------------------------------------
%------------------------------- PREAMBLE --------------------------------------
%-------------------------------------------------------------------------------

%-------------------------------------------------------------------------------
% 	colors
%-------------------------------------------------------------------------------
\usepackage{xcolor}

% define uni-bonn colors
\definecolor{UBBlue}{cmyk}{1.0,0.7,0,0}
\definecolor{UBYellow}{cmyk}{0,0.3,1.0,0}
\definecolor{UBGray}{cmyk}{0,0,0.15,0.55}
\definecolor{UBBlueRGB}{RGB}{7,82,154}
\definecolor{UBYellowRGB}{RGB}{255,185,12}
\definecolor{UBGrayRGB}{RGB}{144,144,133}

% define a certain color style, giving the colors for the background, title, blocks, etc.
\definecolorstyle{uni-bonn-grey}{}
{
	\colorlet{backgroundcolor}{white}
	\colorlet{blocktitlebgcolor}{UBGrayRGB}
	\colorlet{blocktitlefgcolor}{black}
	\colorlet{blockbodybgcolor}{white}
	\colorlet{blockbodyfgcolor}{black}
}

% define another block style, giving the colors for the block title and body
% for blue blocks
\definecolorstyle{uni-bonn-blue}{}{
	\colorlet{blocktitlebgcolor}{UBBlueRGB}
}
% for yellow blocks
\definecolorstyle{uni-bonn-yellow}{}{
	\colorlet{blocktitlebgcolor}{UBYellowRGB}
}

% define a title style, giving the colors for the title and the logo
\definetitlestyle{uni-bonn}{
	width=.8\textwidth,
	linewidth=0pt,
	innersep=8mm,
	titletotopverticalspace=90mm,
	titletoblockverticalspace=20mm,
}{
	% logo figure
	\begin{scope}
		%
		% uni bonn and mctc logo
		\node[anchor=north east, inner sep=0] (uni_logo) at (412.0mm,584.5mm) {\includegraphics[height=80mm]{figures/logo_uni-bonn_mctc.png}};
		%
		% comment the following two block(s) if you dont have 
		% a secondary logo or personal qr code
		%
		% secondary logo 
		\node[anchor=north west, inner sep=0] (imprs_logo) at (-321.0mm,575.0mm) {\includegraphics[height=70mm]{figures/logo_imprs_recharge.png}};
		%
		% personal qr code
		\node[anchor=north west, inner sep=0] (qrcode) at (-415.0mm,575.0mm) {\includegraphics[height=90mm]{figures/qrcode-TF.png}};
		%
		% blue line
		\draw[anchor=north east,
			rounded corners=20pt,
			line width = 6.0mm,
			color=UBBlueRGB]
		(uni_logo.north east) ++(-160mm, -3.3mm) -| ++( -661mm, -1153mm);
		%
		% uncomment the following block if you dont have a secondary logo 
		% and want the gray line from the uni bonn logo
		%
		% gray line
		% \draw[anchor=north east,
		% 	rounded corners=50pt,
		% 	line width = 6.0mm,
		% 	color=UBGrayRGB]
		% (uni_logo.north east) ++(-160mm,-50.95mm) -| ++( -644.4mm, -220mm);
		%
		% yellow line
		\draw[anchor=north east,
			rounded corners=20pt,
			line width = 6.0mm,
			color=UBYellowRGB]
		(uni_logo.north east) ++( -3.07mm,-79mm) |- ++( -805mm, -1091mm);
	\end{scope}
}


% summarize this in a layout theme
\definelayouttheme{uni-bonn}{
	\usecolorstyle{uni-bonn-grey}
	\usebackgroundstyle{Default}
	\usetitlestyle{uni-bonn}
	\useblockstyle{Default}
	\useinnerblockstyle{Default}
	\usenotestyle{Default}
}

% finally, use the theme
\usetheme{uni-bonn}


%-------------------------------------------------------------------------------
% 	fancy layout options
%-------------------------------------------------------------------------------
% \usepackage{layouts}          % for printing page width:
% textwidth in inches: \printinunitsof{in}\prntlen{\textwidth}


%-------------------------------------------------------------------------------
% 	language and time
%-------------------------------------------------------------------------------
\usepackage[T1]{fontenc}
\usepackage[utf8]{inputenc}\DeclareUnicodeCharacter{2212}{-}
\usepackage[american]{babel}
\usepackage[ddmmyyyy]{datetime}
\usepackage{bookmark}

%-------------------------------------------------------------------------------
% 	font
%-------------------------------------------------------------------------------
\usepackage{lmodern}			% use lmodern font
\renewcommand{\familydefault}{\sfdefault}
\linespread{1.2}


%-------------------------------------------------------------------------------
% 	tables
%-------------------------------------------------------------------------------
\usepackage{tabularx}					% table with flexible column width
\usepackage{booktabs}                   % horizontal lines
\usepackage{float} 						% table placement
\usepackage{multirow} 					% multirow table

%-------------------------------------------------------------------------------
% 	figure, graphics
%-------------------------------------------------------------------------------
\usepackage[%
	singlelinecheck=true,
	justification=justified,
	format=plain,
	font=small,
	labelfont=it
]{caption}
\usepackage{subcaption}
\usepackage{graphicx}
\graphicspath{{figures/}}

\usepackage{tikz}

%-------------------------------------------------------------------------------
% 	math
%-------------------------------------------------------------------------------
\usepackage{amsmath,amsfonts,amsthm,amssymb,amstext}
\usepackage[]{sansmath}                   % for math fonts in sans serif
\usepackage{siunitx}
\sisetup{%
detect-all,                             % alle äußeren Schriftkonventionen übernehmen
per-mode=fraction,                      % \per creates fraction
output-decimal-marker={.},              % decimal symbol 
unit-mode=text,                         % Einheiten im Textmodus
separate-uncertainty=true,              % von/bis oder \pm haben beide Einheit nach Zahl
range-phrase=\ {\text{\textendash}}\ ,  % range denoted with endash (note spacing)
range-units=single,                     % range has only one unit
exponent-product=\cdot,                 % use \cdot: \SI{6,626E-34}{\joule\second}
}
\DeclareSIUnit{\calorie}{cal}
\DeclareSIUnit{\ppm}{ppm}
\DeclareSIUnit{\kcal}{\kilo\calorie}


%-------------------------------------------------------------------------------
% 	bibliography, citation, links
%-------------------------------------------------------------------------------
\usepackage[style=english]{csquotes} % required by biblatex.sty
\usepackage[backend=biber,style=chem-angew,maxbibnames=5,sorting=none]{biblatex}
\addbibresource{refs.bib}

\usepackage{hyperref}
\hypersetup{
	colorlinks=true,
	linkcolor=UBBlueRGB,
	citecolor=UBBlueRGB,
	filecolor=UBBlueRGB,
	urlcolor=UBBlueRGB
}
\usepackage[nameinlink]{cleveref}



%-------------------------------------------------------------------------------
%   miscellaneous
%-------------------------------------------------------------------------------
\newcommand{\boltz}{k_{\mathrm{B}}}             % Boltzmann constant
\usepackage{lipsum}

\tikzposterlatexaffectionproofoff               % remove tikz poster watermark


%-------------------------------------------------------------------------------
%------------------------------- END OF PREAMBLE -------------------------------
%-------------------------------------------------------------------------------